% Appendix B (writer: steep-a). Proofs for sec:absorb. Source: steeper notes-final §1.4 (Thm 1), §1.5 (Prop 1.3), §1.7 (Prop 1.7), §4 (Prop 4.1).
\section{Proofs for the absorption section}\label{app:absorb}

\subsection{Classes with an absorbing member}\label{app:absorb:examples}

We prove \cref{rem:absorb:examples}. Suppose first that the mixture $\xi_1=\sum_\pi2^{-\ell(\pi)}P_\pi$ is itself a member, $\xi_1=P_{\pi_U}$. Then $P_\pi\le2^{\ell(\pi)}\xi_1=2^{\ell(\pi)}U$ for every $\pi$, which is (Dom) with $\kappa_0=0$, and $\xi_1=U$. This covers $\Hcl{w}$ ($\Mw=\sum_i\mm(i)\nu_i$ is an lsc weak function semimeasure, hence occurs in the trimmed enumeration), $\Hcl{\mathrm{SI}}$ ($M$ is an lsc semimeasure) and $\Hcl{\mathrm{chr}}$ ($\MH$ is an lsc chronological semimeasure and, by the representation theorem for chronological semimeasures, the law of a history-oracle function machine at every nonempty history; at the empty history both sides of (Dom) are at most $1$ and $P_{\pi_U}=1$) \citep{claudeSFI2026fundamentals}. For $\Hcl{F}$: $\MF$ is a machine law, so it has an index; let $\pi_U$ be a shortest program for it. For $D\ne\emptyset$, $P_\pi(D)=\nu_{T_{V(\pi)}}(D)\le\MF(D)/\mm(V(\pi))\le2^{\ell(\pi)}\MF(D)$, and at $\emptyset$ both sides equal $1$. Here $\xi_1=\MF$ on nonempty data sets and $\xi_1(\emptyset)\le1=\MF(\emptyset)$, so $\xi_1\le U$. \qed

\subsection{Proof of the absorption theorem}\label{app:absorb:absorption}

We prove \cref{thm:absorb:absorption}. Throughout, (Dom) holds and $\lambda>1$.

\emph{(a).} Let $\hat\pi$ be a $\lambda$-MAP for $o$. Then $J_\lambda(\hat\pi;o)\le J_\lambda(\pi_U;o)=\lambda\ell_U-\log U(o)$. By (Dom), $-\log U(o)\le\ell(\hat\pi)+\kappa_0-\log P_{\hat\pi}(o)$. Substituting,
\[
\lambda\ell(\hat\pi)-\log P_{\hat\pi}(o)\le\lambda\ell_U+\ell(\hat\pi)+\kappa_0-\log P_{\hat\pi}(o),
\]
so $(\lambda-1)\ell(\hat\pi)\le\lambda\ell_U+\kappa_0$, which is the length bound. (If $P_{\hat\pi}(o)=0$ then $J_\lambda(\hat\pi;o)=\infty$, and $\hat\pi$ is a minimiser only if every program has infinite objective; we exclude this case, which does not occur when $U(o)>0$.) The loss bound is $J_\lambda(\hat\pi;o)\le J_\lambda(\pi_U;o)$ rearranged. Finally, $\ell^*_\lambda<L$ if and only if $\lambda(L-\ell_U)>L+\kappa_0$.

\emph{(b).} The first inequality keeps the single term of $\pi_U$ in $\xi_\lambda(o)=\sum_\pi2^{-\lambda\ell(\pi)}P_\pi(o)$. The second holds because $2^{-\lambda\ell}\le2^{-\ell}$. If $\xi_1\le2^{\kappa_1}U$, then $-\log\xi_\lambda(o)\ge-\log\xi_1(o)\ge-\log U(o)-\kappa_1$, and the first inequality gives $-\log\xi_\lambda(o)\le\lambda\ell_U-\log U(o)$. For a data-set class, $\xi_\lambda$ is a weak function semimeasure with $\xi_\lambda(\emptyset)\le\sum_\pi2^{-\lambda\ell(\pi)}\le1$, because each $P_\pi$ is one and the weak inequalities are linear. Its sequential log loss on $D_n$ is $-\log\xi_\lambda(D_n)+\log\xi_\lambda(\emptyset)$ in every order (\cref{def:setup:shared}), and the same holds for $U$.

\emph{(c1).} Since $\pi_t$ is deterministic and consistent with $D_n$, $P_{\pi_t}(D_n)=1$, so $\xi_\lambda(D_n)\ge2^{-\lambda\ell_t}$. Also, by (Dom), $U(D_n)\ge2^{-\ell_t-\kappa_0}P_{\pi_t}(D_n)=2^{-\ell_t-\kappa_0}$, so $\xi_\lambda(D_n)\ge2^{-\lambda\ell_U}U(D_n)\ge2^{-\lambda\ell_U-\kappa_0-\ell_t}$. By telescoping, $\sum_{i\le n}-\log\xi_\lambda(t(x_i)\mid x_i,D_{i-1})=-\log\xi_\lambda(D_n)+\log\xi_\lambda(\emptyset)\le-\log\xi_\lambda(D_n)$. Each conditional is at most $1$ by the one-point inequality, so each term is nonnegative, and a label with probability below $1/2$ adds more than one bit.

\emph{(c2).} Combine (a) with $-\log U(D_n)\le\ell_t+\kappa_0$.

\emph{(c3).} Let $\mathcal H_{\le\ell^*_\lambda}:=\{\pi:\ell(\pi)\le\ell^*_\lambda\}$. It is finite by Kraft, and by (a) it contains every $\lambda$-MAP for every $D_n$. For $\pi\in\mathcal H_{\le\ell^*_\lambda}$ let $L_\pi(n):=-\log P_\pi(D_n)$. It is nondecreasing in $n$, because weak function semimeasures are monotone, so it has a limit $L_\pi(\infty)\in[0,\infty]$. Let $\mathcal H_{\rm fin}:=\{\pi\in\mathcal H_{\le\ell^*_\lambda}:L_\pi(\infty)<\infty\}$. It contains $\pi_U$, because $\ell_U\le\ell^*_\lambda$ and $L_{\pi_U}(\infty)\le\ell_t+\kappa_0$. For $\pi\in\mathcal H_{\le\ell^*_\lambda}\setminus\mathcal H_{\rm fin}$, $J_\lambda(\pi;D_n)\to\infty$, while $J_\lambda(\pi_U;D_n)\le\lambda\ell_U+\ell_t+\kappa_0$ for every $n$. Since there are finitely many such $\pi$, there is $n_0$ beyond which none of them is a $\lambda$-MAP. For $\pi\in\mathcal H_{\rm fin}$, $\sum_n[L_\pi(n+1)-L_\pi(n)]\le L_\pi(\infty)$, and each increment equals $-\log\big(P_\pi(D_{n+1})/P_\pi(D_n)\big)$, the loss of $\pi$'s prediction of the next label. So $\pi$ gives the next label probability below $1/2$ fewer than $L_\pi(\infty)+1$ times. Taking the union over the finite set $\mathcal H_{\rm fin}$ bounds the number of $n\ge n_0$ at which some $\lambda$-MAP does so.

\emph{(d).} By (Dom), $\xi_\lambda\ge2^{-\lambda\ell_U}U\ge2^{-\lambda\ell_U-\kappa_0-\ell(\pi_{\Prob})}\Prob$ on every string. Hence, for every length $L$ and $\omega_{1:L}$ drawn from $\Prob$,
\[
\Ex_{\Prob}\log\frac{\Prob(\omega_{1:L})}{\xi_\lambda(\omega_{1:L})}\le\lambda\ell_U+\kappa_0+\ell(\pi_{\Prob}).
\]
Write $\xi_\lambda(b\mid h):=\xi_\lambda(hb)/\xi_\lambda(h)$ and expand by the chain rule over bit positions: $\log(\Prob/\xi_\lambda)(\omega_{1:L})=\sum_{j\le L}\log\big(\Prob(\omega_j\mid\omega_{<j})/\xi_\lambda(\omega_j\mid\omega_{<j})\big)-\log\xi_\lambda(\varepsilon)$, and $-\log\xi_\lambda(\varepsilon)\ge0$. At each position the conditional expectation of the summand is $\KL\big(\Prob(\cdot\mid h)\,\|\,\xi_\lambda(\cdot\mid h)\big)$, which is at least $-\log\sum_b\xi_\lambda(b\mid h)\ge0$ because $\xi_\lambda(\cdot\mid h)$ is a sub-probability. With fixed-length blocks the label positions are fixed. At a label position $\Prob(\cdot\mid h)$ is a point mass, because $y_i=t(x_i)$, so the summand equals $-\log\xi_\lambda(\omega_j\mid\omega_{<j})$; the sum of these over the bits of one label block is $-\log\xi_\lambda(y_i\mid x_1y_1\cdots x_i)$. Dropping the nonnegative terms of the input positions gives the claim. This is Solomonoff's bound \citep{solomonoff1978complexity} restricted to the label positions. \qed

\subsection{Proof of per-input pricing}\label{app:absorb:pricing}

We prove \cref{prop:absorb:pricing}.

\emph{(i), the bound.} $\max_f\ge\Ex_{f\sim Q}$. For each $i$,
\[
\Ex_{f\sim Q}\big[-\log\nu(f(x_i)\mid x_i)\big]=H(Q_{x_i})+\sum_yQ_{x_i}(y)\log\frac{Q_{x_i}(y)}{\nu(y\mid x_i)},
\]
and the last sum is at least $-\log\sum_y\nu(y\mid x_i)\ge0$, by Gibbs' inequality with a sub-probability.

\emph{(i), the value.} For $\varepsilon\in(0,1/|Y|)$ let $\Delta_\varepsilon$ be the set of conditional sub-probabilities on $\{x_1,\dots,x_n\}$ with $\nu(y\mid x_i)\ge\varepsilon$ for all $y,i$. It is compact and convex, and $L(\nu,Q):=\Ex_{f\sim Q}\sum_i-\log\nu(f(x_i)\mid x_i)$ is continuous on $\Delta_\varepsilon\times\Delta(F)$, convex in $\nu$ and linear in $Q$. By Sion's minimax theorem \citep{sion1958general}, $\min_{\Delta_\varepsilon}\max_QL=\max_Q\min_{\Delta_\varepsilon}L$. Since $L$ is linear in $Q$, $\max_QL(\nu,Q)=\max_{f\in F}\sum_i-\log\nu(f(x_i)\mid x_i)$. For each $Q$, $\min_{\Delta_\varepsilon}L(\cdot,Q)$ lies between $\sum_iH(Q_{x_i})$, by the bound, and $\sum_iH(Q_{x_i})-n\log(1-|Y|\varepsilon)$, attained up to that error by $\nu=(1-|Y|\varepsilon)Q_x+\varepsilon$, which lies in $\Delta_\varepsilon$ and sums to $1$. The error term is uniform in $Q$. Hence
\[
\max_Q\sum_iH(Q_{x_i})\le\inf_\nu\max_f\sum_i-\log\nu(f(x_i)\mid x_i)\le\max_Q\sum_iH(Q_{x_i})-n\log(1-|Y|\varepsilon),
\]
where the first inequality is the bound, applied to every $\nu$. Let $\varepsilon\to0$.

\emph{The example.} For $F=\{(a,a),(a,b),(b,a),(c,a)\}$ a direct computation (code in the accompanying repository, seed 2) gives $\max_Q\sum_iH(Q_{x_i})=2.5431$ and the same value for $\min_\nu\max_f$ by constrained optimisation over $\nu$; an independent computation by grid search, mirror ascent and a primal certificate agrees. The naive bound is $\log3+\log2=2.5850$.

\emph{(ii).} For $x\ne0$ exactly half of the parities give each label at $x$, so the uniform $Q$ has $H(Q_x)=1$, and (i) gives a parity with loss at least $n$. The Bayes mixture with weight $2^{-d}$ on each parity gives every parity's data probability at least $2^{-d}$.

\emph{(iii).} Each term of the per-input mixture assigns $1/2$ to the observed label $1$, at every input; the shared mixture gives the constant-$1$ function weight $1/2$. \qed

\subsection{Proof of the price of private randomness}\label{app:absorb:price}

We prove \cref{prop:absorb:price}.

Write $a_i:=0^i$ and $b_i:=1^i$, so $D_\beta=\{(a_i,\beta_i),(b_i,\beta_i):i\le k\}$.

\emph{(i).} Fix $\pi\in\Hcl{\times}$ with conditional semimeasure $\nu=\nu_\pi$. Then $\nu(D_\beta)=\prod_{i\le k}s_i(\beta_i)$ with $s_i(y):=\nu(y\mid a_i)\,\nu(y\mid b_i)$. By the inequality of arithmetic and geometric means, $s_i(y)\le q_i(y)^2$ with $q_i(y):=\frac12\big(\nu(y\mid a_i)+\nu(y\mid b_i)\big)$, and each $q_i$ is a sub-probability on $\{0,1\}$. So $\nu(D_\beta)\le Q_\pi(\beta)^2$ with $Q_\pi(\beta):=\prod_iq_i(\beta_i)\le1$. Since $\sum_\beta Q_\pi(\beta)=\prod_i\sum_yq_i(y)\le1$, $Q_\pi$ is a semimeasure on $\{0,1\}^k$, lower semicomputable uniformly in $(\pi,k)$. Let $\tilde R(\beta):=\sum_{\pi\in\Hcl{\times}}2^{-\ell(\pi)}Q_\pi(\beta)$. It is an lsc semimeasure on $\{0,1\}^k$ given $k$, so $\tilde R(\beta)\le2^{-K(\beta\mid k)+\kappa}$ by (F2). Now let $\gamma:=\min(\lambda,2)$. If $1\le\lambda\le2$, then $2^{-\lambda\ell(\pi)}Q_\pi^2=(2^{-\ell(\pi)}Q_\pi)^\lambda Q_\pi^{2-\lambda}\le(2^{-\ell(\pi)}Q_\pi)^\lambda$. If $\lambda\ge2$, then $2^{-\lambda\ell(\pi)}Q_\pi^2\le(2^{-\ell(\pi)}Q_\pi)^2$. For nonnegative $t_\pi$ and $\gamma\ge1$, $\sum_\pi t_\pi^\gamma\le(\sum_\pi t_\pi)^\gamma$. Therefore
\[
\xi^\times_\lambda(D_\beta)=\sum_\pi2^{-\lambda\ell(\pi)}\nu_\pi(D_\beta)\le\sum_\pi2^{-\lambda\ell(\pi)}Q_\pi(\beta)^2\le\tilde R(\beta)^{\gamma}\le2^{-\gamma(K(\beta\mid k)-\kappa)}.
\]

\emph{(ii).} $2^{-J_\lambda(\pi;D_\beta)}$ is one term of $\xi^\times_\lambda(D_\beta)$. For a product class with a prefix-free code, define $\tilde R$ over that class; (F2) then holds with a constant for the class.

\emph{(iii).} $\nulen$ is the law of the function machine $T=\{(x,w,w_{|x|}):x\ne\varepsilon,\ w\in\{0,1\}^{|x|}\}$, where $w_{|x|}$ is the last bit of the tape prefix $w$. Inputs $a_i,b_i$ have the same length $i$, so they receive the same bit, and distinct $i$ receive independent bits; so $\nulen(D_\beta)=2^{-k}$. Universality gives $\MF\ge2^{-K(\nulen)}\nulen$ \citep{claudeSFI2026fundamentals}, so $-\log\MF(D_\beta)\le k+\kappa$. For $\xi^F_\lambda$ keep the single term of a shortest program for $\nulen$.

\emph{(iv).} Fewer than $2^{k-j}$ strings $\beta$ have $K(\beta\mid k)<k-j$, so $\Pr[K(\beta\mid k)<k-j]<2^{-j}$. On the complement, (i) gives $-\log\xi^\times_\lambda(D_\beta)\ge\gamma(k-j-\kappa)$, while (iii) gives $-\log\MF(D_\beta)\le k+\kappa$ and $-\log\xi^F_\lambda(D_\beta)\le k+\lambda\kappa$. Subtracting, and using $\gamma\le2$ and $\gamma\le\lambda$, the excess is at least $(\gamma-1)k-2j-(\lambda+2)\kappa$ in both cases. \qed

\emph{Computation.} A direct computation (code in the accompanying repository, seed 6) verified the key inequality $\xi^\times_\lambda(D_\beta)\le\tilde R(\beta)^{\min(\lambda,2)}$ in all 24{,}732 cases on random finite product classes with Kraft weights and $\lambda\in\{1,1.3,1.7,2,2.5,4\}$; equality occurs.

\emph{Reading.} At $\lambda=1$ the first option of the product class, writing the coins into the program, costs the same as the shared class up to lower-order terms, which is why $\Mprod$ and $\MF$ are close; at $\lambda>1$ the product class pays $\min(\lambda,2)$ bits per coin, against $1$ for the shared class.

\subsection{Proofs for the two history propositions}\label{app:absorb:history}

We prove \cref{prop:absorb:history}. Each $Z_i$ is a positive real, computable uniformly in the block's history: its terms $w_\theta\prod_{k<i}\nu_\theta(y_k\mid\cdot)$ are computable and positive, and the tail $\sum_{\theta>T}w_\theta\prod_{k<i}\nu_\theta(\cdot)$ is at most $\sum_{\theta>T}w_\theta$, which is computable from $T$ because $w$ and $\sum_\theta w_\theta$ are computable. So $h_w=Z_{i+1}/Z_i$ is computable. It sums to $1$ over $y_i$, because the $\nu_\theta$ are measures. So $h_w$ is a computable conditional measure on extended inputs; a program for it is obtained from a description of $w$ and $(\nu_\theta)$, which gives the length bound in the canonical class. The product of the predictive probabilities over a block of length $n'$ telescopes to $Z_{n'+1}/Z_1$, and $Z_1=\sum_\theta w_\theta\le1$, so
\[
-\log\frac{Z_{n'+1}}{Z_1}\le-\log\frac{w_\theta\prod_i\nu_\theta(y_i\mid\cdot)}{\sum_{\theta'}w_{\theta'}}\le\log\frac1{w_\theta}+\sum_i-\log\nu_\theta(y_i\mid\cdot)
\]
for every $\theta$. \qed

Without the computability assumptions, $Z_{i+1}/Z_i$ is a ratio of lsc sums, which need not be lsc; this is why they are made. A finite family with rational values and rational weights also satisfies them. A direct computation (code in the accompanying repository, seed 5) checked the block identity to $3.6\cdot10^{-15}$ on 2000 random instances with history-dependent inputs.

We prove \cref{prop:absorb:inputs}. A hypothesis of the first kind is an lsc conditional semimeasure on the set of extended inputs, with private randomness per extended input, so it is a product hypothesis there. Extended inputs $(x_{<i},x_i)$ at different positions are distinct. In \cref{prop:absorb:pricing}(ii) take $F$ to be the functions $(x_{<i},x_i)\mapsto\langle a,x_i\rangle\bmod2$ with nonzero $x_i$; under uniform $a$ each label is a fair coin, and the argument is unchanged, as is that of \cref{cor:absorb:nouniversal}. The second claim is \cref{rem:absorb:examples} for $\Hcl{\mathrm{chr}}$. \qed
